% AUTO-GENERATED by writeup/make_rehaul_ablations.py -- do not edit; re-run the script.
% tab:rh_ablation_beyond (RULE.txt addendum G; requested by the author after the HAR + RM + FOMC results).
\begin{table}[htbp]
\centering
\scriptsize
\setlength{\tabcolsep}{2pt}
\begin{threeparttable}
\caption{Beyond HAR, realized moments and the FOMC calendar: change in QLIKE (\%)}
\label{tab:rh_ablation_beyond}
\begin{tabular*}{\linewidth}{@{\extracolsep{\fill}}lrrrrrr@{}}
\toprule
 & \multicolumn{3}{c}{Every bar} & \multicolumn{3}{c}{Ordinary days} \\
\cmidrule(lr){2-4}\cmidrule(l){5-7}
 & Linear & LightGBM & XGBoost & Linear & LightGBM & XGBoost \\
\midrule
\multicolumn{7}{@{}l}{\textit{A. HAR + RM + FOMC plus the group, against HAR + RM + FOMC ($-$: the group adds accuracy)}} \\
Equal-weighted market & $-$0.5 to $-$0.2 & $-$1.0 ($-$4.5) & $-$0.7 ($-$3.3) & $-$0.5 to $-$0.1 & $-$1.0 ($-$5.2) & $-$0.8 ($-$4.3) \\
Value-weighted market & $-$0.8 to $-$0.3 & $-$1.2 ($-$4.6) & $-$1.1 ($-$4.0) & $-$0.7 to $-$0.2 & $-$1.2 ($-$5.0) & $-$1.2 ($-$4.7) \\
Liquidity & $-$0.6 to 0.0 & $-$1.1 ($-$4.4) & $-$1.1 ($-$4.4) & $-$0.6 to $-$0.4 & $-$1.2 ($-$4.7) & $-$1.2 ($-$4.7) \\
Implied volatility & $-$0.3 to 0.3 & $-$0.4 ($-$3.5) & $-$0.5 ($-$4.4) & $-$0.2 to 0.3 & $-$0.5 ($-$3.4) & $-$0.5 ($-$4.8) \\
Sentiment & 0.0 to 0.3 & $-$0.5 ($-$4.7) & $-$0.6 ($-$6.3) & 0.3 to 0.5 & $-$0.5 ($-$4.7) & $-$0.6 ($-$6.9) \\
Volatility demand & 0.0 to 0.3 & $-$0.3 ($-$4.4) & $-$0.3 ($-$4.6) & 0.2 to 0.3 & $-$0.3 ($-$4.6) & $-$0.4 ($-$5.2) \\
All six (the full set) & $-$0.6 to 0.1 & $-$1.6 ($-$4.9) & $-$1.8 ($-$5.3) & $-$0.6 to 0.2 & $-$1.6 ($-$5.3) & $-$1.8 ($-$5.9) \\
Sum of the six rows [ratio to the full set] & $-$1.7 to 0.5 & $-$4.5 [2.88] & $-$4.4 [2.46] & $-$1.1 to 0.4 & $-$4.7 [2.90] & $-$4.7 [2.63] \\
\addlinespace[4pt]
\multicolumn{7}{@{}l}{\textit{B. Full set against HAR + RM + FOMC plus the group ($-$: the full set is more accurate)}} \\
Equal-weighted market & $-$0.2 to 0.2 & $-$0.6 ($-$3.1) & $-$1.1 ($-$4.9) & $-$0.1 to 0.2 & $-$0.6 ($-$3.2) & $-$1.0 ($-$4.6) \\
Value-weighted market & $-$0.2 to 0.3 & $-$0.4 ($-$2.6) & $-$0.7 ($-$4.3) & 0.1 to 0.3 & $-$0.4 ($-$2.7) & $-$0.6 ($-$3.9) \\
Liquidity & $-$0.3 to 0.1 & $-$0.5 ($-$2.6) & $-$0.6 ($-$3.4) & 0.0 to 0.5 & $-$0.4 ($-$2.7) & $-$0.6 ($-$3.6) \\
Implied volatility & $-$0.8 to $-$0.3 & $-$1.1 ($-$3.5) & $-$1.3 ($-$3.8) & $-$0.7 to 0.2 & $-$1.2 ($-$3.8) & $-$1.3 ($-$4.2) \\
Sentiment & $-$0.8 to $-$0.2 & $-$1.1 ($-$3.7) & $-$1.2 ($-$3.6) & $-$0.9 to $-$0.3 & $-$1.1 ($-$4.0) & $-$1.2 ($-$3.9) \\
Volatility demand & $-$0.8 to $-$0.3 & $-$1.2 ($-$3.9) & $-$1.5 ($-$4.5) & $-$0.8 to $-$0.1 & $-$1.3 ($-$4.2) & $-$1.5 ($-$4.8) \\
\bottomrule
\end{tabular*}
\begin{tablenotes}[flushleft]
\footnotesize
\item \textit{Sign:} each cell is $100\,(L_{\text{row}}/L_{\text{ref}} - 1)$, the percentage change in mean QLIKE of the row's forecast against its reference, with the Diebold--Mariano statistic in parentheses (same sign; daily average QLIKE).  Panel~A, a group added to HAR + RM + FOMC (reference: that base; All six: all six groups added, i.e.\ the full set; Sum: the six rows above added up): negative, the loss falls when the group is added, i.e.\ the group adds accuracy; positive, the group hurts.  Panel~B, the full set against HAR + RM + FOMC plus the group (reference: that set): negative, the full set is more accurate.
\item \textit{Notes:} Base: the HAR benchmark set plus the futures' realized moments (RM) and the FOMC calendar (132 columns); each row of Panel~A adds one remaining group to it, Panel~B sets the full set against that model.  Every model keeps the full set's frozen hyperparameters (FOMC-patched design); the random forest is not run.  Linear: the range over ridge, lasso and elastic net.  Ordinary days exclude the \rhabNEventDays{} scheduled FOMC statement days.  The bracket (trees) is the sum of the six Panel~A changes divided by the full set's change: above one, the groups' gains overlap; below one, they complement each other (the linear full sets gain about nothing against the base, so no ratio is shown for them).
\end{tablenotes}
\end{threeparttable}
\end{table}
